\documentclass[12pt]{article}
\usepackage{amsmath, amssymb, geometry, enumitem}
\geometry{margin=1in}
\setlength{\parindent}{0pt}
\setlength{\parskip}{0.5em}

\title{ECON 101: Macroeconomics \\ Quiz 3}
\author{Instructor: Muhebullah Karimzada}
\date{March 05,  2026}

\begin{document}
\maketitle


\section*{Part A: Multiple Choice Questions (10 marks)}

\textit{Each question is worth 2 marks. Choose the best answer.}

\begin{enumerate}[label=\textbf{\arabic*.}]

\item Which of the following would most likely increase long-run economic growth in Canada?
\begin{enumerate}[label=(\Alph*)]
\item An increase in government consumption spending  
\item A reduction in investment in human capital  
\item Greater technological innovation and productivity growth  
\item A decrease in labour force participation  
\end{enumerate}

\item Labour productivity is defined as:
\begin{enumerate}[label=(\Alph*)]
\item Real GDP divided by labour force participation  
\item Real GDP per hour worked  
\item Nominal GDP per worker  
\item Capital stock divided by employment  
\end{enumerate}

\item In the loanable funds model, an increase in government borrowing due to a budget deficit will:
\begin{enumerate}[label=(\Alph*)]
\item Increase the supply of loanable funds  
\item Reduce the demand for loanable funds  
\item Increase interest rates and crowd out private investment  
\item Lower equilibrium interest rates  
\end{enumerate}
\newpage
\item Financial intermediaries contribute to economic growth primarily because they:
\begin{enumerate}[label=(\Alph*)]
\item Eliminate business cycles  
\item Channel savings efficiently into productive investment  
\item Reduce the need for government regulation  
\item Guarantee profits for savers  
\end{enumerate}

\item Which of the following statements best describes the relationship between saving and investment in a closed economy?
\begin{enumerate}[label=(\Alph*)]
\item Saving always exceeds investment  
\item Investment equals national saving  
\item Government spending equals saving  
\item Investment equals consumption  
\end{enumerate}

\end{enumerate}

\hrule
\vspace{1em}

\section*{Part B: Numerical Problem (10 marks)}

Suppose the following data describe an economy:

\[
Y = C + I + G
\]

where:

\[
Y = 5{,}200 \quad
C = 3{,}300 \quad
G = 1{,}200
\]

All values are measured in billions of dollars.



\begin{enumerate}

\item[(a)] Calculate the level of investment in the economy. (3 marks)

\item[(b)] Compute national saving. Show all steps. (3 marks)

\item[(c)] Suppose government spending increases by \$200 billion while taxes remain unchanged.  
Explain how this policy change affects national saving and the equilibrium interest rate in the loanable funds market. (4 marks)

\end{enumerate}

\hrule
\vspace{1em}

\begin{center}
\textit{End of Quiz}
\end{center}

\end{document}